\documentclass[11pt]{article}

\usepackage[utf8]{inputenc}
\usepackage[T1]{fontenc}
\usepackage{lmodern}
\usepackage[margin=1in]{geometry}
\usepackage{amsmath,amssymb,amsthm}
\usepackage{booktabs}
\usepackage{microtype}
\usepackage[hidelinks]{hyperref}
\usepackage{cleveref}

% Run-in headings italic, never bold.
\makeatletter
\renewcommand\paragraph{\@startsection{paragraph}{4}{\z@}%
  {3.25ex \@plus1ex \@minus.2ex}{-1em}%
  {\normalfont\normalsize\itshape}}
\makeatother

\title{The Far End of a Schur Bridge Does Not Care How Its Clusters\\
Are Priced, and the Interior Does}
\author{Peter Cotton\thanks{Microprediction.
Email: \texttt{peter.cotton@microprediction.com}. A companion verification
script is \texttt{paper/verify\_test\_space.py} at
\texttt{github.com/microprediction/schur}. The identities and the
singularity example are checked there in exact rational arithmetic; the
Monte Carlo figures are not.}}
\date{September 20, 2026}

\newcommand{\R}{\mathbb{R}}
\newcommand{\ones}{\mathbf{1}}
\DeclareMathOperator{\diag}{diag}
\DeclareMathOperator{\spn}{span}

\newtheorem{proposition}{Proposition}
\newtheorem{corollary}{Corollary}
\newtheorem{remark}{Remark}
\newtheorem{definition}{Definition}

\begin{document}
\maketitle

\begin{abstract}
A two-tier allocation holds a portfolio inside each cluster and then decides
how much capital each cluster gets. The second decision needs a $k\times k$
matrix, and every construction in this family so far builds that matrix from
the same cluster portfolios it holds. It need not. Separating what is held
from what it is priced against gives a family with two spaces rather than one,
and the far end of the bridge is exact for every choice of the pricing space:
the $k$ conditioned directions already sum to $\Sigma^{-1}u$, so any
non-degenerate pairing recovers it. Choosing one representative asset per
cluster, as hierarchical core-orbital allocation does, therefore cannot affect
the far end. The interior is a different matter. Once the two spaces differ the
$k\times k$ matrix is no longer symmetric, and it can pass through singularity
inside the dial: we give a positive-definite $6\times6$ example where it does,
and on estimated covariances it does so in a quarter to a third of markets. Pricing
against the cluster portfolios keeps the matrix a Gram matrix of $\Sigma$ and
removes the possibility.
\end{abstract}

\section{Introduction}

A two-tier allocation makes two decisions. Inside cluster $C$ it holds some
portfolio, and across clusters it decides how much capital each one gets. The
second decision is an allocation problem in $k$ dimensions, and it needs a
$k\times k$ matrix to work with.

Every construction in this family builds that matrix from the cluster
portfolios themselves. Nested clustered optimization \cite{lopezdeprado2020}
uses the covariance of the cluster portfolios; hierarchical risk parity
\cite{lopezdeprado2016} and hierarchical equal risk contribution
\cite{raffinot2018} use variances computed from the same holdings. Written in
the pair form of \cite{cotton2024schur}, all of them set the pricing object
equal to the holding.

Hierarchical core-orbital allocation \cite{bajo2026} does not. It selects one
physical asset per cluster, allocates across those $k$ assets, and only then
redistributes each cluster's capital among its members. What is priced and what
is held are different objects, and the reason given is institutional: the
strategic decision is expressed in instruments a desk can trade.

This note takes that separation seriously and asks what it costs. The answer is
nothing at the far end and something in the interior. \Cref{sec:far} shows that
the far end is exact for every pricing space, because the $k$ conditioned
directions already sum to the global optimum. \Cref{sec:interior} shows that
once the two spaces differ the $k\times k$ matrix loses its symmetry and can
become singular inside the dial. \Cref{sec:hcoa} reads core-orbital allocation
in these terms.

\section{Two spaces}

Let $\Sigma\succ0$ be the covariance of $n$ assets, let $u\in\R^n$ be a
companion vector, and let $C_1,\dots,C_k$ partition the assets. For a cluster
$C$ write $-C$ for its complement and, for $\gamma\in[0,1]$, define the
conditioned pair
\begin{equation}
  Q_C(\gamma)=\Sigma_{CC}-\gamma\,\Sigma_{C,-C}\Sigma_{-C,-C}^{-1}\Sigma_{-C,C},
  \qquad
  b_C(\gamma)=u_C-\gamma\,\Sigma_{C,-C}\Sigma_{-C,-C}^{-1}u_{-C},
  \label{eq:pair}
\end{equation}
and the direction $d_C(\gamma)=Q_C(\gamma)^{-1}b_C(\gamma)$, embedded in $\R^n$
with support $C$. Collect these as the columns of $D_\gamma\in\R^{n\times k}$.

\begin{definition}[Trial and test spaces]
The \emph{trial} matrix is $D_\gamma$: its column span is what the portfolio is
built from. The \emph{test} matrix is any $R\in\R^{n\times k}$: it is what the
trial directions are priced against. The allocation is
\begin{equation}
  w(\gamma,R)\;\propto\;D_\gamma\big(R^\top\Sigma D_\gamma\big)^{-1}R^\top u ,
  \label{eq:family}
\end{equation}
normalized to sum to one, and defined when $R^\top\Sigma D_\gamma$ is
invertible.
\end{definition}

Taking $R=D_\gamma$ recovers the constructions already in the literature. The
$k\times k$ matrix is then $D_\gamma^\top\Sigma D_\gamma$, and \eqref{eq:family}
is the minimum-variance portfolio on the span of the trial directions. At
$\gamma=0$ with $u=\ones$ that is nested clustered optimization. We call this
the \emph{matched} case, and any $R\neq D_\gamma$ an \emph{unmatched} one.

\section{The far end}
\label{sec:far}

\begin{proposition}[The far end is free of the test space]
\label{prop:far}
Suppose $R^\top\Sigma D_1$ is invertible. Then
$w(1,R)\propto\Sigma^{-1}u$, whatever $R$ is.
\end{proposition}

\begin{proof}
Block inversion on the partition $(C,-C)$ gives
$(\Sigma^{-1}u)_C=Q_C(1)^{-1}b_C(1)=d_C(1)$
\cite[App.~A]{cotton2024schur}. The supports $C_1,\dots,C_k$ partition the
index set, so summing the columns of $D_1$ reconstructs the whole vector:
$D_1\ones_k=\Sigma^{-1}u$. Multiplying on the left by $R^\top\Sigma$,
\[
  R^\top\Sigma D_1\ones_k=R^\top\Sigma\Sigma^{-1}u=R^\top u ,
\]
so $\big(R^\top\Sigma D_1\big)^{-1}R^\top u=\ones_k$ and the right side of
\eqref{eq:family} is $D_1\ones_k=\Sigma^{-1}u$.
\end{proof}

The proposition says something slightly surprising about where the work is
done. At $\gamma=1$ the trial directions are not merely a subspace containing
the answer, they are its blocks, and the coefficients that assemble them are all
one. Any invertible pairing recovers those coefficients, so the pricing space
carries no information at the far end at all.

\begin{corollary}[Representative selection does not reach the far end]
\label{cor:rep}
Let $E\in\R^{n\times k}$ have $E_{p_c,c}=1$ for one chosen asset $p_c$ per
cluster and zeros elsewhere. Then $w(1,E)\propto\Sigma^{-1}u$ for every
selection $p_1,\dots,p_k$ for which $E^\top\Sigma D_1$ is invertible.
\end{corollary}

A rule that chooses the representative, whether by medoid, by alignment with the
cluster's first principal component, or by maximal squared correlation with the
rest of the cluster \cite{bajo2026}, therefore decides where the interior of the
bridge runs and nothing about its far end. The certificate checks
\cref{prop:far} for five test spaces, including two different representative
selections and a dense random one, and for a general companion vector as well as
$u=\ones$.

\begin{remark}[Which far ends]
With $u=\ones$ the far end is minimum variance, with $u=\sigma$ maximum
diversification, and with $u=\mu$ the tangency direction; the proof uses no
property of $u$. \Cref{prop:far} needs the conditioning of \eqref{eq:pair} to be
on every other asset. Conditioning instead on one knot per other cluster is
exact under the model of \cite{cotton2026nco}, and \cref{prop:far} then holds
under that model.
\end{remark}

\section{The interior}
\label{sec:interior}

The freedom of \cref{prop:far} does not extend inward, and the reason is
structural rather than statistical.

\begin{proposition}[The matched case has no pole]
\label{prop:nopole}
If $D_\gamma$ has full column rank then $D_\gamma^\top\Sigma D_\gamma$ is
symmetric positive definite. Hence \eqref{eq:family} with $R=D_\gamma$ is
defined at every $\gamma$ for which the pairs \eqref{eq:pair} exist.
\end{proposition}

\begin{proof}
It is a Gram matrix: $x^\top D_\gamma^\top\Sigma D_\gamma x=(D_\gamma x)^\top
\Sigma(D_\gamma x)>0$ for $x\neq0$ since $\Sigma\succ0$ and $D_\gamma x\neq0$.
\end{proof}

When $R\neq D_\gamma$ the matrix $R^\top\Sigma D_\gamma$ is not symmetric and
\cref{prop:nopole} does not apply. Its determinant is a rational function of
$\gamma$ and may vanish in $[0,1]$ on ordinary positive-definite input. Take
\[
  \Sigma=\begin{pmatrix}
    20 & 8 & -5 & 5 & 5 & -6\\
    8 & 36 & -11 & 8 & 7 & -12\\
    -5 & -11 & 33 & 2 & 6 & 16\\
    5 & 8 & 2 & 16 & 3 & 9\\
    5 & 7 & 6 & 3 & 29 & -18\\
    -6 & -12 & 16 & 9 & -18 & 50
  \end{pmatrix},
\]
with clusters $\{1,2,3\}$ and $\{4,5,6\}$, representatives $3$ and $4$, and
$u=\ones$. Its leading principal minors are positive. The determinant of
$E^\top\Sigma D_\gamma$ is positive at $\gamma=\tfrac12$ and negative at
$\gamma=\tfrac35$, so it vanishes between them and the allocation has a pole
strictly inside the dial. On the same input the matched determinant stays
positive throughout. Both facts are certified in exact arithmetic.

\paragraph{How often this bites.} A pole is not a curiosity of a constructed
example. Take twenty-four assets in four clusters of six, drawn from a
one-factor-per-cluster model, estimate the covariance from one hundred and
twenty observations, and scan the dial on a grid of four hundred points. The
determinant of the outer matrix changes sign somewhere in $[0,1]$ in $28\%$ of
markets when the test space is one asset per cluster and in $35\%$ when it is
dense random. In the matched case it never does, as \cref{prop:nopole} requires.

\paragraph{What it costs out of sample.} The same markets, scored by the true
variance of the estimated portfolio in excess of the true optimum, give the
median profile in \cref{tab:oos}. At the far end all four agree, exactly as
\cref{prop:far} says, and they agree on a number well above zero because the
covariance is estimated. Away from it the unmatched choices are not merely worse
on average, they are occasionally enormous, and the excursions sit where the
determinant crosses.

\begin{table}[ht]
\centering
\begin{tabular}{@{}lrrrr@{}}
\toprule
test space & $\gamma=0$ & $\gamma=0.2$ & $\gamma=0.6$ & $\gamma=1$ \\
\midrule
cluster portfolios (matched) & 10.5 & 9.9 & 12.2 & 23.4 \\
one asset per cluster, maximal squared correlation & 37.3 & 57.0 & 112.9 & 23.4 \\
one asset per cluster, chosen at random & 21.8 & 20.5 & 289.3 & 23.4 \\
dense random & 751.4 & 265.5 & 1112.0 & 23.4 \\
\bottomrule
\end{tabular}
\caption{Excess variance over the true optimum, in percent, mean over thirty
estimated markets. The far-end column is identical across rows by
\cref{prop:far}.}
\label{tab:oos}
\end{table}

So the separation of the two spaces buys exactness that is free and stability
that is not. A test space chosen for interpretability is safe at the far end and
exposed everywhere else.

\section{Core-orbital allocation}
\label{sec:hcoa}

Hierarchical core-orbital allocation \cite{bajo2026} selects a core asset $p_c$
in each cluster, allocates capital $W$ across the $k$ cores using their own
covariance $\Sigma_{\mathcal{R}\mathcal{R}}=E^\top\Sigma E$, and then
redistributes each $W_c$ among the members of cluster $c$ in proportion to a
portfolio computed from that cluster's own block. Writing $V$ for the matrix of
those normalized cluster portfolios, the result is $w=VW$.

In the notation of \eqref{eq:family} the trial matrix is $D_0$, whose columns are
the cluster portfolios up to scale, and the test matrix is $E$. The outer
matrix, however, is $E^\top\Sigma E$ and not the pairing $E^\top\Sigma D_0$. The
method is therefore adjacent to the family rather than inside it, and the
difference is one factor, which the rest of this section makes precise.

\begin{proposition}[When the two coincide]
\label{prop:hcoa}
Let $N=\diag(\ones^\top d_{C_c}(0))$, so that $D_0=VN$. Then
$E^\top\Sigma D_0=(E^\top\Sigma V)N$, and core-orbital allocation agrees with
the member $w(0,E)$ of \eqref{eq:family} for every companion vector if and only
if
\begin{equation}
  E^\top\Sigma\,(V-E)=0 ,
  \label{eq:gap}
\end{equation}
that is, if and only if each cluster's holding has the same covariance with
every core as that core has.
\end{proposition}

Condition \eqref{eq:gap} holds trivially when every cluster is a singleton,
where the holding is its own core and core-orbital allocation is the far end
itself. Under a one-factor model within each cluster it says that each holding
carries the same factor exposure as its core. It does not hold in general, and
the certificate exhibits a case where it fails.

\paragraph{Which matrix to price with.} Pricing with $E^\top\Sigma E$ rather
than with $E^\top\Sigma D_0$ places the method outside the family, so
\cref{prop:far} does not apply to it. That is not a defect to be repaired. A
method fixed at $\gamma=0$ never visits the far end, so the exactness it forgoes
is of no use to it, while $E^\top\Sigma E$ is a principal submatrix of $\Sigma$
and therefore symmetric positive definite, which is exactly what
\cref{sec:interior} shows the pairing is not. Substituting the pairing at
$\gamma=0$ on estimated covariances lowers the median excess variance and raises
the worst case by an order of magnitude in our runs, so the substitution is not
one we would recommend.

\paragraph{What the choice does cost.} Under a one-factor model within each
cluster the covariance of the holdings satisfies $V^\top\Sigma V=K\Sigma_{\mathcal{R}
\mathcal{R}}K+\Delta$, with $K$ diagonal in each cluster's factor exposure and
$\Delta$ diagonal, so pricing by the cores omits both. Minimum variance is
invariant to a common scaling, so the omission is immaterial when the clusters
carry similar exposures and grows with their dispersion. On simulated universes
whose clusters resemble one another, as the maturity and credit segments of a
bond curve do, the two price vectors differ by half a percentage point of weight
and their out-of-sample variances by a fraction of a percent; on universes whose
clusters differ sharply in exposure and volatility the gap is several times
larger. The comparability with nested clustered optimization reported in
\cite{bajo2026} on a fixed-income universe is what this predicts, and an equity
universe with dispersed cluster exposures is where the prediction could fail.

\section{Closing}

The bridges catalogued so far are named by their two ends, and each assumes that
a cluster is priced by the portfolio held in it. Dropping that assumption adds a
second axis rather than another row. Along it the far end is not a point but a
face: any non-degenerate pricing space reaches $\Sigma^{-1}u$, so the choice of
representative is a statement about the interior alone.

Two questions are left open. The first is whether any pricing space improves on
the matched one in the interior, where the evidence here is that none of the
three tried does. The second is whether a pricing space can be chosen to avoid
the pole while keeping exactness, for instance by requiring
$R^\top\Sigma D_\gamma$ to be positive definite along the path; the matched
choice satisfies this by construction, and whether any other does is not settled
here.

\begin{thebibliography}{99}

\bibitem{bajo2026}
M. Bajo Traver.
Hierarchical core-orbital allocation: decoupling clustering from optimization in
portfolio construction.
\emph{The Journal of Financial Data Science}, forthcoming, 2026.

\bibitem{cotton2024schur}
P. Cotton.
Schur complementary allocation: a unification of hierarchical risk parity and
minimum variance portfolios.
arXiv:2411.05807, 2024.

\bibitem{cotton2026nco}
P. Cotton.
Nested clustered optimization is one end of a Schur bridge, and the interior is
sometimes provably better.
SSRN 7480738, 2026.

\bibitem{lopezdeprado2016}
M. L\'opez de Prado.
Building diversified portfolios that outperform out of sample.
\emph{Journal of Portfolio Management}, 42(4):59--69, 2016.

\bibitem{lopezdeprado2020}
M. L\'opez de Prado.
\emph{Machine Learning for Asset Managers}.
Cambridge University Press, 2020.

\bibitem{raffinot2018}
T. Raffinot.
The hierarchical equal risk contribution portfolio.
SSRN 3237540, 2018.

\end{thebibliography}

\end{document}
